\chapter{\Acl*{FIBnt}}

\section{Artefacts of the \acs*{FIBnt} process}

Due to the destructive nature of the \ac{FIBnt} process, there are a lot of possible artefacts.

\begin{figure}[H]
    \centering

	\begin{subfigure}[t]{0.32\textwidth}
		\includegraphics[width=\textwidth]{methods/fib_artifacts_b.pdf}
		\caption{Redeposition in front of the \ac{ROI} and slight vertical shadowing.}
		\label{fig:fib_artifacts_b}
	\end{subfigure}
	\hfill
	\begin{subfigure}[t]{0.32\textwidth}
		\includegraphics[width=\textwidth]{methods/fib_artifacts_c.pdf}
		\caption{Charging effects in form of drift artefacts and redepositioning.}
		\label{fig:fib_artifacts_c}
	\end{subfigure}
	\hfill
	\begin{subfigure}[t]{0.32\textwidth}
		\includegraphics[width=\textwidth]{methods/fib_artifacts_a.pdf}
		\caption{Curtaining and charging effects.}
		\label{fig:fib_artifacts_a}
	\end{subfigure}
    \caption{Imaging issues of hydrated and embedded \ac{C3S} specimens during the \ac{FIBnt}-process within the bulk material.}
    \label{fig:fib_artifacts}
\end{figure}

As in the \ac{ArBIB} sectioning (see \zcref{sec:bib_pol_fixed}), the cutting process can introduce curtaining.
It is caused by the cutting process causing vertical line artefacts in the direction of the ion beam (\zcref{fig:fib_artifacts_a}).
These artefacts can be reduced in post processing using a directional \ac{FFT} based filter as for the curtaining found in the \ac{ArBIB} process \cite{fitschen.2017,schankula.2018}.
However, this processing reduces image quality and avoiding the artefact in the first place would be preferable.

Vertical shadowing is another typical artefact caused by a loss of signal of the back-scattered electrons due to the trench geometry (see \zcref{fig:fib_artifacts_b}).
This well known effect can be corrected in post processing \cite{holzer.2007, mac.2018, fager.2020}, if the phase contrast is good enough.
If the signal-to-noise ratio is to low in the darkened areas, this post processing wont help.
However, there are more severe uncorrectable artefacts caused by the in-place slicing process.
Since cementitious material is not conductive and the trenches are not coated throughout, charging effects can negatively influence the image quality (e.g. \zcref{fig:fib_artifacts_c} and \zcref{fig:fib_artifacts_a}) or cause uncorrectable image drift during the process.
The re-deposition of the removed material usually happens within the trenches, which does not cause any problems.
However, it may also appear on small structures in front of the \ac{ROI} or even on the specimen surface after longer imaging periods as shown in \zcref{fig:fib_artifacts_b} and \zcref{fig:fib_artifacts_c}.

\subsection{The lift-out process}

\begin{figure}[h!]
    \centering

	\includegraphics[width=\textwidth]{methods/fib_liftout/image095_colored.pdf}
    \caption{\qty{15}{\micro\meter} wide specimen ({\color{green}A}).
	On the left-hand side is the needle welded to the specimen ({\color{blue}B}). They are navigated and welded to the copper comb ({\color{red}C}).
	The platinum vapour source ({\color{orange}D}) used for welding is visible on the right. \cite{kleiner.2021c}}
    \label{fig:fib_liftout}
\end{figure}

To avoid or at least reduce issues caused by slicing within the bulk material, the \ac{ROI}  can be completely cut free.
Then a lift-out process can be applied, which is similar to the preparation process of \ac{TEM} lamellas \cite{overwijk.1993, langford.2001a}.
To achieve this, a needle is inserted and welded to the \ac{ROI} cube using the platinum coating process.
Afterwards, the cube is completely cut loose from the bulk material and a copper comb is inserted into the SEM.
Finally, the cube is transferred and subsequently welded to a finger of the comb as shown in \zcref{fig:fib_liftout}.

\begin{figure}[htb]
    \centering

    \newlength{\desiredheight}
    \setlength{\desiredheight}{0.219\textheight}


	\begin{subfigure}[t]{0.7\textwidth}
		\includegraphics[height=\desiredheight]{methods/fib_liftout/ion_image.pdf}%
		\caption{Top view (\ac{SEM} image using the ion beam.)}
		\label{fig:fib_liftout_result_a}
	\end{subfigure}
	\hfill
	\begin{subfigure}[t]{0.29\textwidth}
		\includegraphics[height=\desiredheight]{methods/fib_liftout/bse_image.pdf}
		\caption{Front view (\ac{SEM}-\ac{BSE} image)}
		\label{fig:fib_liftout_result_b}
	\end{subfigure}

    \caption{
		The cement cube (\qtyproduct{20 x 20 x 20}{\micro\meter\cubed}) welded to the copper comb. \cite{kleiner.2025c}
	}
    \label{fig:fib_liftout_result}
\end{figure}

The prepared and separated cube can now be serially cut as illustrated in \zcref{fig:fib_scheme} with significantly reduced shadowing correlating to the gel pores and charging effects, respectively.

A detailed video tutorial of the \ac{FIBnt} process, the required parameters, and potential sources of errors for cement-based materials was created and published \cite{kleiner.2025b}.

\FloatBarrier
\section{Evaluation obstacles}
\label{sec:fib_evaluation_obstacles}

While the specimen is vacuum-infused with a low-viscosity resin, it often occurs that some larger pores are not completely filled.
This complicates the segmentation and therefore the analysis of the pore space, since needle-like structures tend to move and the pore background is constantly visible.
Such small structures are therefore ignored in the segmentation process, especially if \ac{EDX} datasets are available, since they allow distinction between pore and phase due to a reduced signal strength.

Due to the sensitive nature of the hydrate phases (see \zcref{sec:beam_damage}), the operator must pay particular attention when analysing at higher voltages as for \ac{EDX} or \ac{EBSD} samples.
As shown in \zcref{fig:edx_damage}, an electron beam at \qty{10}{\kilo\volt}, \qty{4}{\pico\ampere}, and a long exposure time (\qty{600}{\micro\second}, \num{6} frames) can cause damage in the specimen which could be misinterpreted as porosity.
The exposure time is the sum of the dwell time on each point and the frame count.
But especially the dwell time increases the risk of specimen damage.
In the example, pores with a diameter of several \unit{\nano\meter} were induced, which can mostly be identified by comparing the damaged slice (\zcref{fig:edx_damage_a}) with a freshly cut surface (\zcref{fig:edx_damage_b}).
In this case, the damage was visible in a depth of $\approx$ \qty{350}{\nano\meter}.

\begin{figure}[htb]
    \centering

	\begin{subfigure}[t]{0.49\textwidth}
		\includegraphics[width=\textwidth]{methods/edx-damage/a.pdf}%
		\caption{An undamaged surface slice.}
		\label{fig:edx_damage_a}
	\end{subfigure}
	\hfill
	\begin{subfigure}[t]{0.49\textwidth}
		\includegraphics[width=\textwidth]{methods/edx-damage/b.pdf}
		\caption{Damaged surface after \ac{EDX}-measurements.}
		\label{fig:edx_damage_b}
	\end{subfigure}
    \caption{
        Sections of \ac{BSE}-images (\qty{2}{\kilo\volt}, \qty{2}{\pico\ampere}) created during a \ac{FIBnt}-process illustrating the effect of \ac{EDX} at \qty{10}{\kilo\volt}, \qty{4}{\pico\ampere} after \num{6} frames and a dwell time of \qty{600}{\micro\second} on a hydrated \ac{C3S} surface, prepared using \ac{FIB}.
    }
	% details in E:\Mitarbeiter\Florian Kleiner\FIB-Data\C3S + EDX\2020_11_18 C3S 28d+EDX (2 Teile)\EDX_RAW\2020_11_18 C3S 28d+EDX
    \label{fig:edx_damage}
\end{figure}

Solutions to this problem are to reduce energy density.
This can be done by reducing the scanning duration or the electron beam current, which subsequently will reduce the overall counts, the electron beam voltage, which will limit the elements which can be analysed or to use more sensitive detectors.
For the presented experiments, an even more reduced electron beam voltage was not applicable.
Therefore, a very sensitive, windowless \ac{EDX}-detector was used and combined with the counts of a second, less-sensitive \ac{EDX}-detector, which yields only one-fifth of the counts of the windowless detector.
While a reduction of the accelerating voltage from \num{10} to \qty{6}{\kilo\volt} had no significant effect, a reduced dwell time (reduced from \num{600} to \qty{60}{\micro\second}) and increased frame count (from \num{6} to \num{60}) significantly decreased the damage.

While these measures result in a good signal-to-noise ratio for the main components, especially the less abundant elements have a fairly noisy mapping.
However, it still allows the distinction of the different main phases of the material.


\section{Post processing following the lift out process}

Samples prepared via the lift-out process and their resulting datasets require much less post-processing to reduce artefacts.
However, some issues still remain to be solved.
For example, noisy \ac{EDX} datasets have to be denoised with an edge-preserving denoising algorithm.
Since the signal-to-noise ratio of the \ac{EDX} element mappings are very low, algorithms like median filtering and \ac{NLM} are not suitable.
In contrast, \ac{BM3D} and \ac{TV} (using the split \textsc{Bregman} optimisation-based denoising) were found to be effective for the noise found in the \ac{EDX} slices and enable reproducible quantitative analysis across elements.
The parameters employed for \ac{BM3D} were: $h$=40, \texttt{templateWindowSize}=4, \texttt{searchWindowSize}=8, \texttt{groupSize}=32, and \texttt{slidingStep}=3.
For \ac{TV} using the split \textsc{Bregman} method, the parameters were: \texttt{weight}=\num{20.0}, \texttt{max\_num\_iter}=500, $\epsilon$=\num{0.0001}.
The \ac{BM3D} results seemed to be slightly sharper and more consistent, especially for low-signal elements like \ce{S}, when using the same denoising parameters for the whole dataset (see \zcref{fig:edx_noise}).
For example, the \ac{TV} filter introduces circular artefacts in the \ce{S} mapping as shown in \zcref{fig:edx_noise_e}.
However, the processing time for \ac{BM3D} is significantly higher than for the \ac{TV} method.

%D:\2025_01_07 FIB-Zylinder\EDX\Ca\Slice 1151-Ca K_alpha_1EDS - SliceImage - 1151 118.tif
\begin{figure}[htb]
    \centering

	\begin{subfigure}[t]{0.32\textwidth}
		\includegraphics[width=\textwidth]{methods/fib_postprocessing/Ca_noisy.pdf}
		\caption{Noisy \ce{Ca} mapping}
		\label{fig:edx_noise_a}
	\end{subfigure}
	\hfill
	\begin{subfigure}[t]{0.32\textwidth}
		\includegraphics[width=\textwidth]{methods/fib_postprocessing/Ca_TV_B.pdf}
		\caption{Denoised \ce{Ca} mapping (\ac{TV})}
		\label{fig:edx_noise_b}
	\end{subfigure}
	\hfill
	\begin{subfigure}[t]{0.32\textwidth}
		\includegraphics[width=\textwidth]{methods/fib_postprocessing/Ca_BM3D.pdf}
		\caption{Denoised \ce{Ca} mapping (\ac{BM3D})}
		\label{fig:edx_noise_c}
	\end{subfigure}
	\hfill
	\begin{subfigure}[t]{0.32\textwidth}
		\includegraphics[width=\textwidth]{methods/fib_postprocessing/S_noisy.pdf}
		\caption{Noisy \ce{S} mapping}
		\label{fig:edx_noise_d}
	\end{subfigure}
	\hfill
	\begin{subfigure}[t]{0.32\textwidth}
		\includegraphics[width=\textwidth]{methods/fib_postprocessing/S_TV_B.pdf}
		\caption{Denoised \ce{S} mapping (\ac{TV})}
		\label{fig:edx_noise_e}
	\end{subfigure}
	\hfill
	\begin{subfigure}[t]{0.32\textwidth}
		\includegraphics[width=\textwidth]{methods/fib_postprocessing/S_BM3D.pdf}
		\caption{Denoised \ce{S} mapping (\ac{BM3D})}
		\label{fig:edx_noise_f}
	\end{subfigure}
    \caption{
      Effect of denoising using the \ac{TV} (split \textsc{Bregman} optimisation) and \ac{BM3D} filters on the raw\ac{EDX}-mappings of \ce{Ca} (high signal) and \ce{S} (low signal).
    }
    \label{fig:edx_noise}
\end{figure}

The main challenge for many datasets is the precise alignment of the acquired images.
Spatial drift during acquisition can cause a shift of the observed object between consecutive slices.
For a single image stack, this issue can be addressed effectively using the \ac{SIFT} algorithm as implemented in \textit{TrakEM2} \cite{cardona.2012}, which allows an automated correction of in-plane translations.
Furthermore, after the \ac{SIFT} alignment the stack is often still systematically shifted about \qtyrange{30}{40}{\degree} caused by the \ac{SEM} geometry (see \zcref{fig:fib_scheme}).

Aligning the \ac{EDX} dataset to the corresponding \ac{BSE} dataset is challenging.
Although the correspondence between \ac{BSE} and \ac{EDX} slices is known, the images differ in resolution and field of view.
Usually the image shift between the \ac{BSE} and \ac{EDX} images of the same slice should be nearly identical.
This allows to solve the alignment problem relatively easily.
Depending on the data acquisition settings it can, however, happen that the \ac{EDX} dataset does not align perfectly with the \ac{BSE} dataset (e.g. due to drift caused by charging).
Then, additional registration steps, such as the use of landmarks or scale-adjusted transformations, are required to achieve accurate overlay of the datasets.

\subsection{Development of an evaluation workflow}

Therefore, a Python tool utilising \textit{napari} \cite{sofroniew.2025}, a graphical user interface to view and edit 2D and 3D image data, was developed.
This workflow denoises, aligns, visualises, and analyses correlative \ac{BSE} and \ac{EDX} image stacks.
First, the pre-processed in-plane translation data stored in the \textit{TrakEM2} project are loaded.
The user must provide the image scale and slice thickness in \unit{\nano\meter\per\px}, unless this information is available in the file metadata.

Following denoising, inter-slice intensity banding ($z$-banding) in the 3D-\ac{EDX} dataset and brightness fluctuations in the 3D-\ac{BSE} dataset are mitigated by normalisation of each slice, based on the mean brightness value of the entire stack (see \zcref{fig:sem_z-banding}).

\begin{figure}[htb]
	\centering

	\begin{subfigure}[t]{0.49\textwidth}
		\includegraphics[width=\textwidth]{methods/fib_postprocessing/z-band-raw_scaled.pdf}
		\caption{Uncorrected}
		\label{fig:sem_z-banding_a}
	\end{subfigure}
	\hfill
	\begin{subfigure}[t]{0.49\textwidth}
		\includegraphics[width=\textwidth]{methods/fib_postprocessing/z-band-corr_scaled.pdf}
		\caption{Corrected}
		\label{fig:sem_z-banding_b}
	\end{subfigure}
	\caption{
		Illustration of the mean brightness variations between \ac{EDX} slices ($z$-banding) and the result of the correction.
		The images depict the $x$-$z$ plane and the mean values of the denoised \ac{EDX} image stack.
	}
	\label{fig:sem_z-banding}
\end{figure}

Subsequently, at least three (preferably more) landmark pairs between the first \ac{BSE} and \ac{EDX} images must be specified manually.
These landmarks are utilised to estimate an isotropic scale-and-translation transformation, thereby enabling accurate stack-to-stack alignment.


\begin{figure}[htb]
    \centering

	\begin{subfigure}[t]{0.32\textwidth}
		\includegraphics[width=\textwidth]{methods/fib_postprocessing/BSE.pdf}
		\caption{\ac{BSE} image}
		\label{fig:sem_postprocessing_a}
	\end{subfigure}
	\hfill
	\begin{subfigure}[t]{0.32\textwidth}
		\includegraphics[width=\textwidth]{methods/fib_postprocessing/EDX_CaSiAl.pdf}
		\caption{
			\ac{EDX}-map ({\color{red}\ce{Ca}}, {\color{green}\ce{Si}} and {\color{blue}\ce{Al}})
		}
		\label{fig:sem_postprocessing_b}
	\end{subfigure}
	\hfill
	\begin{subfigure}[t]{0.32\textwidth}
		\includegraphics[width=\textwidth]{methods/fib_postprocessing/superpixel.pdf}
		\caption{Superpixel and clustering (in 3D)}
		\label{fig:sem_postprocessing_c}
	\end{subfigure}
    \caption{
      Postprocessing of the dataset illustrated on a single slice of the \ac{FIBnt} volume.
	  The \ac{SLIC}O superpixel segmentation (small multicoloured areas) is half transparently overlaid with an automatic clustering of these superpixels (larger similar coloured areas). \cite{kleiner.2025b}
    }
    \label{fig:sem_postprocessing}
\end{figure}

Since the electron beam of the \ac{SEM} hits the \ac{ROI} under an angle of \qty{38}{\degree}, the image stacks inherently exhibit a systematic vertical shift following alignment with the \ac{SIFT} algorithm.
This shift can be compensated by applying a user-defined angle for the correction.
This required angle is primarily determined by the device geometry and the orientation at which the specimen cube is welded to the copper comb.
The image distortion caused by this geometry effect can and should be corrected for during the recording of the dataset.

The co-registered and corrected \ac{BSE} and \ac{EDX} voxel datasets are visualised using physically meaningful scales.
The user can select a \ac{ROI}, which is then utilised for the subsequent processing.
Within the \ac{ROI}, three-dimensional \ac{SLIC}O superpixels are computed on the denoised \ac{EDX} volume (see \zcref{fig:sem_postprocessing_c}).
In analogy to the approach of Multi\ac{SLIC} \cite{yu.2023}, three principal elements (\ce{Ca}, \ce{Si}, \ce{Al}) are chosen.
This allows to assign specific elements to individual (red-green-blue) colour channels in an image, from which the superpixels are subsequently generated.
This enables the use of standard \ac{SLIC}O implementations implemented in the \textit{OpenCV} \cite{bradski.2000} library.
For each superpixel, the mean intensity for each element is calculated from the raw \ac{EDX} data.
These superpixel-based vectors are subsequently normalised and clustered using $k$-means in combination with \ac{UMAP}.
This approach facilitates the identification of material phases, resulting in a three-dimensional phase label volume that is directly overlaid and visualised in \textit{napari} \cite{sofroniew.2025}.

% TODO: compare different ways to get to the goal? - SLIC -> (direct/indirect with umap?) -> labeling; UNET
% TODO: how good is the data alignment? What does this to the segmentation - see ANtons dataset
% Wie ist das mit der Reliefstruktur und der Abschattung während EDX?





\section{The laser preparation process}

Using \ac{FIBnt} alone, practically limits the size of the \ac{ROI} due to the machine time required.
Furthermore, issues such as redeposition, charging and shadowing are primarily caused by geometric constraints of the \ac{FIBnt} process, including the presence of narrow trenches.
While it is possible to enlarge the trenches using \ac{FIB}, doing so is often impractical and time-consuming.
Laser ablation offers an alternative for rapid, high-volume trench creation.
An example of such a prepared specimen is shown in \zcref{fig:laser_prep_tilt_anim}.

\begin{figure}[htbp]
    \centering

    \begin{subfigure}[t]{0.49\textwidth}
		\ifdefined\PRINTABLE
			\includegraphics[width=\textwidth]{methods/laser_prep/tilt6.pdf}
		\else
			\animategraphics[palindrome,poster=6,autoplay,width=\textwidth]{5}{methods/laser_prep/tilt}{0}{6}
		\fi
        \caption{Laser-prepared bar and trenches.}
        \label{fig:laser_prep_tilt_anim}
    \end{subfigure}
    \hfill
    \begin{subfigure}[t]{0.49\textwidth}
        \includegraphics[width=\textwidth]{methods/laser_prep/C3S_7d_Laser_prep_069.pdf}
        \caption{Fresh \ac{FIBnt} slice.}
        \label{fig:fibnt_fresh_slice}
    \end{subfigure}
    \caption{
		\qty{7}{\d} hydrated \ac{C3S} specimen, embedded in epoxy resin and prepared using a \ac{fs-laser}.
		The preparation site has a size of \qtyproduct{.6 x 1.0}{\milli\meter\squared} and the prepared bar is \qty{40}{\micro\meter} wide at the top.
		The depth of the removed trenches is estimated to be more than \qty{100}{\micro\meter}.
        The white deposition on the top is a protective \ce{Pt} coating.
    }
    \label{fig:laser_prep}
\end{figure}

The cross-section (\zcref{fig:fibnt_fresh_slice}) shows minor damage to the bulk material.
Amorphous depositions up to \qty{2}{\micro\meter} and debris-filled pores are visible on the surface sectioned using a \ac{FIB}.
However, specimens which were not lifted out of the bulk material cannot be analysed using \ac{EBSD} due to geometrical limitations.

% TODO volume


The laser preparation also facilitates the easy preparation of larger cross-sections of \ac{ROI}s.
For example, it enables preparation of mortar samples to analyse the \ac{ITZ} between binder and aggregate.
\zcref{fig:laser_prep_itz} shows that the cement binder is ablated much faster than the aggregates.

\begin{figure}[htbp]
    \centering

	\includegraphics[width=\textwidth]{methods/laser_prep/fib-zylinder.pdf}
    \caption{
		\ac{SEM}-\ac{SE} image of a tilted (\qty{52}{\degree}) mortar sample prepared using a \ac{fs-laser}.
    	The higher ablation rate of cement binder ($\approx\qty{165}{\micro\meter}$) compared to the quartz aggregate ($\approx\qty{45}{\micro\meter}$) is clearly visible.
		In the rectangle ({\color{darkred}red}), the area of the \ac{FIB} preparation site is highlighted.
		The \ac{ITZ} is marked with the dotted line.
    }
    \label{fig:laser_prep_itz}
\end{figure}

Nevertheless, the laser prepared trenches enable the creation of very large \ac{FIB}-polished areas of the \ac{ITZ}, which also allows acquisition of \ac{EDX} datasets (\zcref{fig:laser_prep_itz_fib}).

\begin{figure}[htbp]
    \centering

    \begin{subfigure}[t]{0.49\textwidth}
		\includegraphics[width=\textwidth]{methods/laser_prep/itz_fib_edx_a.pdf}

        \caption{
			\ac{BSE} image of a section with few artefacts.
		}
        \label{fig:laser_prep_itz_fib_a}
    \end{subfigure}
    \hfill
    \begin{subfigure}[t]{0.49\textwidth}
        \includegraphics[width=\textwidth]{methods/laser_prep/itz_fib_edx_b.pdf}
        \caption{
			\ce{Ca}-\ac{EDX} mapping overlaid on top of the \ac{BSE} of a later slice.
		}
        \label{fig:laser_prep_itz_fib_b}
    \end{subfigure}
    \caption{
		\ac{FIB} polished section of an \ac{ITZ} of a mortar.
		The binder matrix can be found in both images on the left side, the quartz aggregate on the right.
		The \ce{Ca}-rich \ac{ITZ} is marked with dotted white lines.
    }
    \label{fig:laser_prep_itz_fib}
\end{figure}

However, due to the orientation of the beam, the large-scale surface polishing using \ac{FIB} still introduces typical artefacts.
For example, protruding edges cast shadows on the \ac{BSE} image (\zcref{fig:laser_prep_itz_fib}, left; note the shadow on the homogeneous grey value of the aggregate on the right side).
Furthermore, the softer cementitious binder tends to be more prone for curtaining artefacts, while the quartz grain appears to be nicely polished.
Finally, since a very large area of an non-conductive material is exposed to the electron beam, charging effects occur even on low voltage and current combinations (\qty{5}{\kilo\volt}, \qty{0.1}{\nano\ampere}) with low dwell times (\qty{50}{\nano\second}).
Therefore, the reachable chemical contrast in the \ac{BSE} images is very limited.

The \ac{EDX} mapping also suffers from significant signal attenuation with increasing distance from the surface.
For example, the \ac{EDX} mapping \zcref{fig:laser_prep_itz_fib} has been brightened in the lower portions to compensate for this effect, which also increases the noise in these areas.
Nevertheless, a zone of higher \ce{Ca} concentrations at the interface to the quartz grain becomes visible.
The width of this layer ranges from \qtyrange{2}{5}{\micro\meter}, which is smaller, than the expected \ac{ITZ} width (\qtyrange{20}{80}{\micro\meter} \cite{lavrov.2017,torsaeter.2015,zhu.2000}).
Therefore, this \ce{Ca}-rich zone, which most likely consists of \ac{CH} as shown in \zcref{fig:laser_prep_itz}, is only a part of the \ac{ITZ}.
